\section{AHSME 1980}
        \begin{enumerate}
        	\item (\textit{AHSME 1980}) The largest whole number such that seven times the number is less than $100$ is


 $\text{(A)} \ 12 \qquad \text{(B)} \ 13 \qquad \text{(C)} \ 14 \qquad \text{(D)} \ 15 \qquad \text{(E)} \ 16$


	\item (\textit{AHSME 1980}) The degree of $(x^2+1)^4 (x^3+1)^3$ as a polynomial in $x$ is


 $\text{(A)} \ 5 \qquad \text{(B)} \ 7 \qquad \text{(C)} \ 12 \qquad \text{(D)} \ 17 \qquad \text{(E)} \ 72$


	\item (\textit{AHSME 1980}) If the ratio of $2x-y$ to $x+y$ is $\frac{2}{3}$, what is the ratio of $x$ to $y$?


 $\text{(A)} \ \frac{1}{5} \qquad \text{(B)} \ \frac{4}{5} \qquad \text{(C)} \ 1 \qquad \text{(D)} \ \frac{6}{5} \qquad \text{(E)} \ \frac{5}{4}$


	\item (\textit{AHSME 1980}) In the adjoining figure, $CDE$ is an equilateral triangle and $ABCD$ and $DEFG$ are squares. The measure of $\angle GDA$ is


 $\text{(A)} \ 90^\circ \qquad \text{(B)} \ 105^\circ \qquad \text{(C)} \ 120^\circ \qquad \text{(D)} \ 135^\circ \qquad \text{(E)} \ 150^\circ$

                    \begin{figure}[H]
                        \centering
                        \includegraphics[width=0.5\textwidth]{Problem/AHSME_I_1980_4_Problem_Figure_1.png}
                    \end{figure}
                    


	\item (\textit{AHSME 1980}) If $AB$ and $CD$ are perpendicular diameters of circle $Q$, and $\measuredangle QPC = 60^\circ$, then the length of $PQ$ divided by the length of $AQ$ is 


 
                    \begin{figure}[H]
                        \centering
                        \includegraphics[width=0.5\textwidth]{Problem/AHSME_I_1980_5_Problem_Figure_1.png}
                    \end{figure}
                    
$\text{(A)} \ \frac{\sqrt{3}}{2} \qquad \text{(B)} \ \frac{\sqrt{3}}{3} \qquad \text{(C)} \ \frac{\sqrt{2}}{2} \qquad \text{(D)} \ \frac12 \qquad \text{(E)} \ \frac23$


	\item (\textit{AHSME 1980}) A positive number $x$ satisfies the inequality $\sqrt{x} < 2x$ if and only if


 $\text{(A)} \ x > \frac{1}{4} \qquad \text{(B)} \ x > 2 \qquad \text{(C)} \ x > 4 \qquad \text{(D)} \ x < \frac{1}{4}\qquad \text{(E)} \ x < 4$


	\item (\textit{AHSME 1980}) Sides $AB,BC,CD$ and $DA$ of convex polygon $ABCD$ have lengths $3$, $4$, $12$, and $13$, respectively, and $\measuredangle CBA$ is a right angle. The area of the quadrilateral is


 
                    \begin{figure}[H]
                        \centering
                        \includegraphics[width=0.5\textwidth]{Problem/AHSME_I_1980_7_Problem_Figure_1.png}
                    \end{figure}
                    
$\text{(A)} \ 32 \qquad \text{(B)} \ 36 \qquad \text{(C)} \ 39 \qquad \text{(D)} \ 42 \qquad \text{(E)} \ 48$


	\item (\textit{AHSME 1980}) How many pairs $(a,b)$ of non-zero real numbers satisfy the equation


 
 \[\frac{1}{a} + \frac{1}{b} = \frac{1}{a+b}\]
$\text{(A)} \ \text{none} \qquad \text{(B)} \ 1 \qquad \text{(C)} \ 2 \qquad \text{(D)} \ \text{one pair for each} ~b \neq 0$
$\text{(E)} \ \text{two pairs for each} ~b \neq 0$


	\item (\textit{AHSME 1980}) A man walks $x$ miles due west, turns $150^\circ$ to his left and walks $3$ miles in the new direction. If he finishes a a point $\sqrt{3}$ from his starting point, then $x$ is


 $\text{(A)} \ \sqrt 3 \qquad \text{(B)} \ 2\sqrt{5} \qquad \text{(C)} \ \frac 32 \qquad \text{(D)} \ 3 \qquad \text{(E)} \ \text{not uniquely determined}$


	\item (\textit{AHSME 1980}) The number of teeth in three meshed circular gears $A$, $B$, and $C$ are $x$, $y$, and $z$, respectively. (The teeth on all gears are the same size and regularly spaced.) The angular speeds, in revolutions per minute, of $A$, $B$, and $C$ are in the proportion


 $\text{(A)} \ x: y: z ~~\text{(B)} \ z: y: x ~~ \text{(C)} \ y: z: x~~ \text{(D)} \ yz: xz: xy ~~ \text{(E)} \ xz: yx: zy$


	\item (\textit{AHSME 1980}) If the sum of the first $10$ terms and the sum of the first $100$ terms of a given arithmetic progression are $100$ and $10$, 
respectively, then the sum of first $110$ terms is


 $\text{(A)} \ 90 \qquad \text{(B)} \ -90 \qquad \text{(C)} \ 110 \qquad \text{(D)} \ -110 \qquad \text{(E)} \ -100$


	\item (\textit{AHSME 1980}) The equations of $L_1$ and $L_2$ are $y=mx$ and $y=nx$, respectively. Suppose $L_1$ makes twice as large of an angle with the horizontal (measured counterclockwise from the positive x-axis ) as does $L_2$, and that $L_1$ has 4 times the slope of $L_2$. If $L_1$ is not horizontal, then $mn$ is


 $\text{(A)} \ \frac{\sqrt{2}}{2} \qquad \text{(B)} \ -\frac{\sqrt{2}}{2} \qquad \text{(C)} \ 2 \qquad \text{(D)} \ -2 \qquad \text{(E)} \ \text{not uniquely determined}$


	\item (\textit{AHSME 1980}) A bug (of negligible size) starts at the origin on the coordinate plane. First, it moves $1$ unit right to $(1,0)$. Then it makes a $90^\circ$ counterclockwise and travels $\frac 12$ a unit to $\left(1, \frac 12 \right)$. If it continues in this fashion, each time making a $90^\circ$ degree turn counterclockwise and traveling half as far as in the previous move, to which of the following points will it come closest? 


 $\text{(A)} \ \left(\frac 23, \frac 23 \right) \qquad \text{(B)} \ \left( \frac 45, \frac 25 \right) \qquad \text{(C)} \ \left( \frac 23, \frac 45 \right) \qquad \text{(D)} \ \left(\frac 23, \frac 13 \right) \qquad \text{(E)} \ \left(\frac 25, \frac 45 \right)$


	\item (\textit{AHSME 1980}) If the function $f$ defined by 

 \[f(x)=\frac{cx}{2x+3} ,\quad x\neq -\frac{3}{2} ,\] 
satisfies $x=f(f(x))$ for all real numbers $x$ except $-\frac{3}{2}$, then $c$ is


 $\text{(A)} \ -3 \qquad  \text{(B)} \ - \frac{3}{2} \qquad  \text{(C)} \ \frac{3}{2} \qquad  \text{(D)} \ 3 \qquad  \text{(E)} \ \text{not uniquely determined}$


	\item (\textit{AHSME 1980}) A store prices an item in dollars and cents so that when $4\%$ sales tax is added, no rounding is necessary because the result is exactly $n$ dollars where $n$ is a positive integer. The smallest value of $n$ is


 $\text{(A)} \ 1 \qquad \text{(B)} \ 13 \qquad \text{(C)} \ 25 \qquad \text{(D)} \ 26 \qquad \text{(E)} \ 100$


	\item (\textit{AHSME 1980}) Four of the eight vertices of a cube are the vertices of a regular tetrahedron. Find the ratio of the surface area of the cube to the surface area of the tetrahedron.


 
                    \begin{figure}[H]
                        \centering
                        \includegraphics[width=0.5\textwidth]{Problem/AHSME_I_1980_16_Problem_Figure_1.png}
                    \end{figure}
                    
$\text{(A)} \ \sqrt 2 \qquad \text{(B)} \ \sqrt 3 \qquad \text{(C)} \ \sqrt{\frac{3}{2}} \qquad \text{(D)} \ \frac{2}{\sqrt{3}} \qquad \text{(E)} \ 2$


	\item (\textit{AHSME 1980}) Given that $i^2=-1$, for how many integers $n$ is $(n+i)^4$ an integer?


 $\text{(A)} \ \text{none} \qquad \text{(B)} \ 1 \qquad \text{(C)} \ 2 \qquad \text{(D)} \ 3 \qquad \text{(E)} \ 4$


	\item (\textit{AHSME 1980}) If $b>1$, $\sin x>0$, $\cos x>0$, and $\log_b \sin x = a$, then $\log_b \cos x$ equals


 $\text{(A)} \ 2\log_b(1-b^{a/2}) ~~\text{(B)} \ \sqrt{1-a^2} ~~\text{(C)} \ b^{a^2} ~~\text{(D)} \ \frac 12 \log_b(1-b^{2a}) ~~\text{(E)} \ \text{none of these}$


	\item (\textit{AHSME 1980}) Let $C_1, C_2$ and $C_3$ be three parallel chords of a circle on the same side of the center. 
The distance between $C_1$ and $C_2$ is the same as the distance between $C_2$ and $C_3$. 
The lengths of the chords are $20, 16$, and $8$. The radius of the circle is 


 $\text{(A)} \ 12 \qquad  \text{(B)} \ 4\sqrt{7} \qquad  \text{(C)} \ \frac{5\sqrt{65}}{3} \qquad  \text{(D)}\ \frac{5\sqrt{22}}{2}\qquad \text{(E)}\ \text{not uniquely determined}$


	\item (\textit{AHSME 1980}) A box contains $2$ pennies, $4$ nickels, and $6$ dimes. Six coins are drawn without replacement, 
with each coin having an equal probability of being chosen. What is the probability that the value of coins drawn is at least $50$ cents? 


 $\text{(A)} \ \frac{37}{924} \qquad  \text{(B)} \ \frac{91}{924} \qquad  \text{(C)} \ \frac{127}{924} \qquad  \text{(D)}\ \frac{132}{924}\qquad \text{(E)}\ \text{none of these}$


	\item (\textit{AHSME 1980}) 
                    \begin{figure}[H]
                        \centering
                        \includegraphics[width=0.5\textwidth]{Problem/AHSME_I_1980_21_Problem_Figure_1.png}
                    \end{figure}
                    


 In triangle $ABC$, $\measuredangle CBA=72^\circ$, $E$ is the midpoint of side $AC$, 
and $D$ is a point on side $BC$ such that $2BD=DC$; $AD$ and $BE$ intersect at $F$. 
The ratio of the area of $\triangle BDF$ to the area of quadrilateral $FDCE$ is


 $\text{(A)} \ \frac 15 \qquad  \text{(B)} \ \frac 14 \qquad  \text{(C)} \ \frac 13 \qquad  \text{(D)}\ \frac{2}{5}\qquad \text{(E)}\ \text{none of these}$


	\item (\textit{AHSME 1980}) For each real number $x$, let $f(x)$ be the minimum of the numbers $4x+1, x+2$, and $-2x+4$. Then the maximum value of $f(x)$ is 


 $\text{(A)} \ \frac{1}{3} \qquad  \text{(B)} \ \frac{1}{2} \qquad  \text{(C)} \ \frac{2}{3} \qquad  \text{(D)} \ \frac{5}{2} \qquad  \text{(E)}\ \frac{8}{3}$


	\item (\textit{AHSME 1980}) Line segments drawn from the vertex opposite the hypotenuse of a right triangle to the points trisecting the hypotenuse have lengths $\sin x$ and $\cos x$, where $x$ is a real number such that $0<x<\frac{\pi}{2}$. The length of the hypotenuse is 


 $\text{(A)} \ \frac{4}{3} \qquad  \text{(B)} \ \frac{3}{2} \qquad  \text{(C)} \ \frac{3\sqrt{5}}{5} \qquad  \text{(D)}\ \frac{2\sqrt{5}}{3}\qquad \text{(E)}\ \text{not uniquely determined}$


	\item (\textit{AHSME 1980}) For some real number $r$, the polynomial $8x^3-4x^2-42x+45$ is divisible by $(x-r)^2$. Which of the following numbers is closest to $r$? 


 $\text{(A)} \ 1.22 \qquad  \text{(B)} \ 1.32 \qquad  \text{(C)} \ 1.42 \qquad  \text{(D)} \ 1.52 \qquad  \text{(E)} \ 1.62$


	\item (\textit{AHSME 1980}) In the non-decreasing sequence of odd integers $\{a_1,a_2,a_3,\ldots \}=\{1,3,3,3,5,5,5,5,5,\ldots \}$ each odd positive integer $k$
appears $k$ times. It is a fact that there are integers $b, c$, and $d$ such that for all positive integers $n$, 
$a_n=b\lfloor \sqrt{n+c} \rfloor +d$, 
where $\lfloor x \rfloor$ denotes the largest integer not exceeding $x$. The sum $b+c+d$ equals 


 $\text{(A)} \ 0 \qquad  \text{(B)} \ 1 \qquad  \text{(C)} \ 2 \qquad  \text{(D)} \ 3 \qquad  \text{(E)} \ 4$


	\item (\textit{AHSME 1980}) Four balls of radius $1$ are mutually tangent, three resting on the floor and the fourth resting on the others. 
A tetrahedron, each of whose edges have length $s$, is circumscribed around the balls. Then $s$ equals 


 $\text{(A)} \ 4\sqrt 2 \qquad  \text{(B)} \ 4\sqrt 3 \qquad  \text{(C)} \ 2\sqrt 6 \qquad  \text{(D)}\ 1+2\sqrt 6\qquad \text{(E)}\ 2+2\sqrt 6$


	\item (\textit{AHSME 1980}) The sum $\sqrt[3] {5+2\sqrt{13}}+\sqrt[3]{5-2\sqrt{13}}$ equals 


 $\text{(A)} \ \frac 32 \qquad  \text{(B)} \ \frac{\sqrt[3]{65}}{4} \qquad  \text{(C)} \ \frac{1+\sqrt[6]{13}}{2} \qquad  \text{(D)}\ \sqrt[3]{2}\qquad \text{(E)}\ \text{none of these}$


	\item (\textit{AHSME 1980}) The polynomial $x^{2n}+1+(x+1)^{2n}$ is not divisible by $x^2+x+1$ if $n$ equals 


 $\text{(A)} \ 17 \qquad  \text{(B)} \ 20 \qquad  \text{(C)} \ 21 \qquad  \text{(D)} \ 64 \qquad  \text{(E)} \ 65$


	\item (\textit{AHSME 1980}) How many ordered triples $(x,y,z)$ of integers satisfy the system of equations below? 


 
 \[\begin{array}{l} x^2-3xy+2y^2-z^2=31 \\ -x^2+6yz+2z^2=44 \\ x^2+xy+8z^2=100\\ \end{array}\]
$\text{(A)} \ 0 \qquad  \text{(B)} \ 1 \qquad  \text{(C)} \ 2 \qquad  \text{(D)}\ \text{a finite number greater than 2}\qquad \text{(E)}\ \text{infinitely many}$


	\item (\textit{AHSME 1980}) A six digit number (base $10$) is squarish if it satisfies the following conditions: 


 (i) none of its digits are zero; 


 (ii) it is a perfect square; and 


 (iii) the first of two digits, the middle two digits and the last two digits of the number are all perfect squares when considered as two digit numbers. 


 How many squarish numbers are there? 


 $\text{(A)} \ 0 \qquad  \text{(B)} \ 2 \qquad  \text{(C)} \ 3 \qquad  \text{(D)} \ 8 \qquad  \text{(E)} \ 9$


\end{enumerate}